\section{WebGL 기반 디지털 플랫폼 기술 --- 서비스 구조 및 제작 요구사항}

본 장은 three.js 기반 WebGL 3D 렌더링이 ARES에서 어떻게 활용되는지를
랜딩 페이지, 미션 시뮬레이터, 오프라인 빌드 측면에서 분석하고, 웹앱 기반
제작 시 요구사항을 정리한다.

\subsection{three.js 번들링 전략}

three.js는 CDN이 아닌 로컬 사전 번들로 동봉된다.

\begin{itemize}
  \item \code{Web/vendor/three-bundle.min.js} --- \code{window.THREE} 네임스페이스와,
        \code{GLTFLoader}\textbf{·}\code{OrbitControls}\textbf{·}\code{RoomEnvironment}를
        담은 \code{window.ARES3} 네임스페이스를 노출한다.
  \item \code{Web/vendor/meshopt\_decoder.js} --- GLB 모델이
        \code{EXT\_meshopt\_compression}으로 압축되어 있어, \code{GLTFLoader}가
        파싱하기 전에 반드시 먼저 로드되어야 한다.
\end{itemize}

로컬 번들 채택 이유는 \code{file://}\textbf{·}오프라인 환경 지원이다. CDN
페치 의존을 제거함으로써 인터넷 없는 교실에서도 3D 렌더링이 동작한다.

\subsection{랜딩 페이지 3D (\code{index.html})}

랜딩 페이지는 로버 모델을 인터랙티브하게 렌더링한다.

\begin{itemize}
  \item \textbf{렌더러}: \code{WebGLRenderer}(antialias, alpha), sRGB 출력,
        ACES 필름릭 톤 매핑, 소프트 섀도우(\code{PCFSoftShadowMap}).
  \item \textbf{환경광}: \code{RoomEnvironment} 기반 PMREM 환경맵 + 3점 조명
        (반구광\textbf{·}키 라이트\textbf{·}필 라이트).
  \item \textbf{카메라/컨트롤}: \code{PerspectiveCamera}(약 40\textdegree\ FOV),
        \code{OrbitControls}로 좌우 회전 허용, 상하 틸트 $\pm$10\textdegree\ 제한,
        줌\textbf{·}팬 비활성.
  \item \textbf{본 애니메이션}: 모델의 팔 본(\code{RightArm},
        \code{RightForeArm})을 쿼터니언 보간으로 움직여 ``손 흔들기'' 인사를 구현.
\end{itemize}

\subsubsection{\code{file://}에서의 모델 로딩}

랜딩 로버 모델은 \code{Mesh/\allowbreak ares\_robot.glb}와, 이를 base64로
임베드한 \code{Mesh/\allowbreak ares\_robot.embed.js}
(\code{window.\allowbreak ARES\_ROBOT\_GLB\_B64})로 제공된다.
프로토콜에 따라 분기한다.

\captionof{listing}{file:// 분기 로딩}
\begin{minted}{javascript}
if (location.protocol === 'file:') {
  ensureEmbed()                       // base64 임베드 스크립트 로드
    .then(b64 => loader.parse(b64ToBuffer(b64), '', onModelLoad, onModelErr));
} else {
  loader.load('Mesh/ares_robot.glb', onModelLoad, undefined, onModelErr);
}
\end{minted}

\code{file://}에서는 \code{fetch} 대신 base64 ArrayBuffer를 직접
\code{GLTFLoader.parse()}에 넘겨 오프라인 로딩을 보장한다.

\subsection{미션 시뮬레이터 (\code{simulation.js})}

\code{simulation.js}는 실제 BLE 연결 없이 블록 프로그램을 3D로 시연하는
시뮬레이터이다. 주제(\code{TOPICS})별로 다른 환경 모델과 연출을 로드한다.

\begin{longtable}{L{2.8cm} L{4.2cm} L{6.8cm}}
\toprule
\rowcolor{tablehead}\textbf{주제} & \textbf{모델} & \textbf{연출} \\
\midrule
\endhead
알비 & \code{Mesh/AlbiStaticLow.glb} & 눈 LED 구체\textbf{·}점광, OLED 캔버스 텍스처 \\
신호등 & \code{Mesh/LampBox.glb} 외 & 신호등 램프\textbf{·}손 모델 \\
발사대 & \code{Mesh/LaunchStation.glb} & 안테나/로켓 메시 분리 재질, 레이더, 발사 LED 스트립 \\
로버 & \code{Mesh/RoverParts/*.glb} & 7개 부품(본체\textbf{·}건\textbf{·}헤드\textbf{·}LED\textbf{·}OLED\textbf{·}레이더\textbf{·}휠) 조립 \\
\bottomrule
\end{longtable}

시뮬레이터는 블록 실행 시 명령을 UI 로그로 출력하고, 눈 LED\textbf{·}OLED 텍스처를
실시간 갱신하며, 발사대 주제에서는 지오메트리 인덱스를 분석해 안테나/로켓을
별도 메시로 분리하는 후처리(\code{recolorLaunchpadAntenna})를 수행한다.

\subsection{오프라인 단독 실행본 빌드}

\code{build.ps1}(Windows) / \code{build.sh}(macOS\textbf{·}Linux) /
\code{build.bat}(진입점)는 \code{Web/}을 \code{file://}로 동작하는 단일
\code{Build/} 폴더로 변환한다. 주요 단계는 다음과 같다.

\begin{enumerate}
  \item \code{Build/} 초기화.
  \item \textbf{ES 모듈 번들링}: \code{esbuild}로 \code{main.js}와 의존 모듈을
        단일 IIFE \code{main.bundle.js}로 묶음.
  \item \textbf{외부 라이브러리 로컬화}: Blockly 압축본을 \code{vendor/}로 내려받음.
  \item \textbf{정적 자산 복사}: 대시보드, CSS, three.js 번들, meshopt 디코더,
        임베드 로봇 등.
  \item \textbf{\code{index.html}\textbf{·}\code{main.html} 패치}:
        \code{type="module"} 제거, \code{defer} 스크립트 체인으로 전환,
        경로 평탄화.
  \item \textbf{자산 인라인화}:
    \begin{itemize}
      \item GLB 바이너리는 모델별 \code{vendor/bin\_<이름>.js}로 분할
            (base64; GitHub 100\,MB 파일 제한 회피).
      \item 텍스트 자산(\code{overview.html}, \code{lesson.json}, 예제 XML)은
            \code{vendor/inline\_assets.js}의 \code{DATA}에 등록.
      \item \code{window.fetch}를 가로채는 심을 설치하여, \code{Mesh/*.glb}\textbf{·}
            텍스트 요청을 인라인 데이터로 응답.
    \end{itemize}
\end{enumerate}

\captionof{listing}{inline\_assets.js의 fetch 심(개념)}
\begin{minted}{javascript}
window.fetch = function (input, init) {
  const url = norm(input);
  const bk = find(BIN, url);   // 바이너리(GLB) 우선 매칭
  if (bk !== null)
    return Promise.resolve(new Response(b64ToU8(BIN[bk]),
      { status: 200, headers: { 'Content-Type': 'application/octet-stream' } }));
  const tk = find(DATA, url);  // 텍스트 자산
  if (tk !== null)
    return Promise.resolve(new Response(DATA[tk],
      { status: 200, headers: { 'Content-Type': mimeFor(tk) } }));
  return origFetch ? origFetch(input, init) : Promise.reject(...);
};
\end{minted}

\code{defer} 로딩 순서는 (1) \code{bin\_*.js}가 \code{window.\_\_ARES\_BIN\_\_}을
채우고, (2) \code{inline\_assets.js}가 fetch 심을 설치하며, (3)
\code{main.bundle.js}가 앱을 시작하도록 보장된다. 이 순서와 런타임 자산
가로채기의 연계는 그림~\ref{fig:buildload}와 같다. 핵심은 앱이 평소처럼
\code{fetch('Mesh/...glb')}를 호출하더라도, 설치된 심이 이를 가로채 인라인된
base64를 \code{Response}로 돌려주므로 \emph{앱 코드를 수정하지 않고도} 오프라인
동작이 가능하다는 점이다.

\begin{diagram}{오프라인 빌드 자산 로딩 \textbf{·} fetch 심 연계도}{fig:buildload}
\node[boxw, text width=24mm] (bin) {① bin\_*.js\\{\scriptsize defer}};
\node[boxw, below=7mm of bin, text width=24mm] (inl) {② inline\_assets.js\\{\scriptsize defer}};
\node[boxw, below=7mm of inl, text width=24mm] (bun) {③ main.bundle.js\\{\scriptsize 앱 시작}};
\node[store, right=14mm of bin, text width=24mm] (store) {\_\_ARES\_BIN\_\_\\{\scriptsize GLB base64}};
\node[box, right=14mm of inl, text width=26mm] (shim) {window.fetch 심\\{\scriptsize BIN/DATA 조회}};
\node[boxw, right=14mm of bun, text width=26mm] (call) {fetch('Mesh/x.glb')};
\node[hw, right=12mm of shim, text width=26mm] (resp) {Response\\{\scriptsize base64$\to$Uint8Array}};
\node[hw, above=7mm of resp, text width=26mm] (gltf) {GLTFLoader\\$\to$ 3D 장면};
\draw[flow] (bin)--(inl); \draw[flow] (inl)--(bun);
\draw[flow] (bin) -- node[lbl,above]{채움} (store);
\draw[flow] (inl) -- node[lbl,above]{설치} (shim);
\draw[flow] (bun) -- (call);
\draw[flow] (call.north) to[bend right=12] node[lbl,right]{가로채기} (shim.south);
\draw[dep] (store) -- (shim);
\draw[flow] (shim) -- (resp);
\draw[flow] (resp) -- (gltf);
\end{diagram}

\subsection{웹앱 기반 제작 시 요구사항}

\begin{infobox}[title=WebGL 자산 제작 \textbf{·} 배포 체크리스트]
\begin{itemize}
  \item GLB 모델은 \code{EXT\_meshopt\_compression}으로 압축하고, 디코더를
        \code{GLTFLoader}보다 먼저 로드한다.
  \item \code{file://} 지원을 위해 모든 모델 로딩에 \code{fetch} 또는
        \code{parse} 경로를 일관되게 두고, 오프라인 빌드에서 fetch 심으로
        가로챌 수 있게 한다.
  \item 대용량 GLB는 모델별 base64 청크로 분할하여 단일 파일 크기 제한을 피한다.
  \item three.js\textbf{·}Blockly 등은 로컬 사본으로 동봉하여 CDN 비의존성을 확보한다.
  \item 스크립트 로딩 순서를 \code{defer} 체인으로 명시하여 전역 네임스페이스
        준비 순서를 보장한다.
\end{itemize}
\end{infobox}

\begin{teachbox}{오프라인으로 만드는 법}
\teachhead{설계 의도}
\begin{itemize}
  \item three.js\textbf{·}Blockly를 로컬 동봉하고, \code{fetch} 심으로 자산을 인라인해 \code{file://} 단독 실행을 완성한다.
  \item \code{defer} 로딩 순서가 동작의 전제다.
\end{itemize}
\teachhead{분석할 때 핵심 질문}
\begin{itemize}
  \item \code{fetch} 심이 없으면 \code{file://}에서 무엇이 실패하는가?
  \item 왜 GLB를 모델별 base64 청크로 나눴는가?
\end{itemize}
\teachhead{점검 요소}
\begin{itemize}
  \item 빌드 스크립트(\code{build.sh}/\code{build.ps1})를 직접 실행해 \code{Build/}를 만들고, 브라우저로 더블클릭해 동작을 확인하라.
\end{itemize}
\end{teachbox}

